\chapter{My Idea}
\label{chap:my_idea}

So, my idea is to create a simple google chrome extension that will look
something like this:

\begin{figure}[h]
    \centering
    \incfig{outline}
    \caption{Here's how the extension will look like}
    \label{fig:outline}
\end{figure}

\section{The Outline Explained}
\label{sec:the_outline_explained}

In this section, I will explain the outline of the extension in details
including how everything will work.

\subsection{File Browser}
\label{sub_sec:file_browser}

This section's used to list all of your notes. The first tab (\textbf{notes})
will list all of your notes. The second tab (\textbf{lessons}) will list all of
the lessons within the current course.

You can also click and drag the tabs around. For example, you can click and drag
the \textbf{Lessons} tab to the bottom of the \textbf{File Browser}, then it
will make both of them show. One at the top and one at the bottom.

% subsection file_browser (end)

\subsection{Main Section}
\label{sub_sec:main_section}

You click the \textbf{notes} tab to go to the current lesson's notes. The
\textbf{Compiled PDF} tab will just show you the compiled PDF.
If you need help with \textbf{LaTeX}, then you click the \textbf{Help} tab. It
will open another tab and bring you to \textbf{overleaf.com/learn}. That's a
great website that teaches you how to use \textbf{LaTeX} from beginner to pro.
The final tab is the \textbf{Review} tab. This is probebly my favorite tab.
Once you click it, it will list all of the notes that you have created in the
current course. After you select one of them, it will start to ask you review
questions about the lesson's content. But, unfortunately, you have to create
those questions manually when you are taking notes. But, eventually, I want to
create an AI bot that does that manually.

% subsection main_section (end)

\subsection{History Watcher}
\label{sub_sec:history_watcher}

This is the section that monitors the history of your notes. If you know
anything about \textbf{Git}, then this section is very self-explanatory.

Here's how it works:

When you modify a file, it will show up in the history.

After you've finished modifing the file, there's going to be a button on the
bottom that you click to save the modified file to the history.

Here's the format:

\textbf{Date of Modification (Time of Modification) - Author of Modification}

On the buttom of the modification line, there will be a button that if you
click, it will expand a section to show you all of the modified files.

% subsection history_watcher (end)

\subsection{Output Box}
\label{sub_sec:output_box}

This is the box that shows you the errors and warnings that the
\textbf{LaTeX Compiler} gives you.
You can drag the border around the section to make it bigger or smaller.

% subsection output_box (end)

% section the_outline_explained (end)

% chapter my_idea (end)

\newpage
